\section{Các bối cảnh triển khai tại bệnh viện}
\label{sec:contexts}

Tiền đề của dự án này là không có một kiến trúc duy nhất nào phù hợp với mọi
bệnh viện. Mục này đặc trưng hoá các bối cảnh tồn tại trong hệ thống y tế Việt
Nam theo những chiều thực sự quyết định điều gì có thể triển khai được.

\subsection{Cấu trúc hệ thống bệnh viện Việt Nam}

Y tế công lập tại Việt Nam được tổ chức theo tuyến, từ các bệnh viện tuyến
trung ương trực thuộc Bộ Y tế, qua các bệnh viện tuyến tỉnh và tuyến huyện, tới
các trạm y tế xã. Năng lực chẩn đoán hình ảnh bám theo phân tuyến này: CT là
thường quy từ tuyến tỉnh trở lên và ngày càng phổ biến ở tuyến huyện, trong khi
MRI tập trung ở tuyến trung ương và các cơ sở tuyến tỉnh lớn. Nhân lực bác sĩ
chẩn đoán hình ảnh phân bố mất cân đối hơn nhiều so với trang thiết bị, và
chính sự bất đối xứng này, có máy chụp mà thiếu người đọc, xác định cơ hội cho
một nền tảng cho phép hình ảnh đến được với bác sĩ ở bất kỳ đâu.

\subsection{Các chiều phân biệt}

\begin{description}[leftmargin=3.0cm,style=nextline]
    \item[Hạ tầng hình ảnh] Có hay không hệ thống PACS, hệ thống đó có hỗ trợ
          DICOMweb hay không, và hình ảnh được lưu trữ tập trung hay được giữ
          lại trên máy trạm của thiết bị chụp.
    \item[Tài nguyên tính toán] Có sẵn phần cứng GPU hay không, có hạ tầng máy
          chủ hay không, và có thể giả định về an ninh vật lý cũng như tính
          liên tục của nguồn điện hay không.
    \item[Kết nối mạng] Băng thông và độ tin cậy của đường truyền ra mạng bên
          ngoài, yếu tố quyết định liệu một thành phần đặt từ xa có khả thi hay
          không.
    \item[Năng lực CNTT] Có hay không nhân sự đủ khả năng vận hành, bảo trì hệ
          thống đã triển khai, yếu tố quyết định mức độ phức tạp vận hành có
          thể chấp nhận.
    \item[Nhân lực lâm sàng] Có bác sĩ chẩn đoán hình ảnh tại chỗ hay không và
          trong khung giờ nào, yếu tố quyết định mức độ cần thiết của khả năng
          truy cập hình ảnh từ xa.
    \item[Khối lượng] Số ca chụp mỗi ngày, quyết định yêu cầu về thông lượng.
\end{description}

\subsection{Các mẫu bối cảnh điển hình}

Ba mẫu bối cảnh điển hình được đề xuất, định nghĩa sao cho một bệnh viện cụ thể
có thể được xếp vào một mẫu chỉ bằng quan sát.

\begin{table}[H]
\centering
\small
\begin{tabularx}{\textwidth}{@{}l X X X@{}}
\toprule
 & \textbf{A: Trung ương} & \textbf{B: Tuyến tỉnh} & \textbf{C: Tuyến huyện} \\
\midrule
PACS & Hoàn thiện, DICOMweb & Có, mức độ khác nhau & Thường không có \\
Lưu trữ ảnh hiện tại & PACS/VNA, có sao lưu & PACS dung lượng hạn chế, sao lưu không đều & Ổ cứng máy trạm thiết bị, bị ghi đè khi đầy, không sao lưu \\
GPU & Có sẵn & Có thể có & Không \\
Kết nối mạng & Ổn định, tốc độ cao & Trung bình & Không liên tục \\
Nhân sự CNTT & Phòng ban chuyên trách & 1--2 người đa năng & Không có tại chỗ \\
Bác sĩ chẩn đoán hình ảnh & Tại chỗ, 24/7 & Giờ hành chính & Không có hoặc luân phiên \\
Khối lượng (CT sọ não/ngày) & $>50$ & 10--50 & $<10$ \\
Nhu cầu chính & Thông lượng, tích hợp PACS sẵn có & Lưu trữ và truy cập ngoài giờ & Lưu trữ tại chỗ, kết nối tuyến trên \\
\bottomrule
\end{tabularx}
\caption{Các mẫu bối cảnh bệnh viện điển hình.}
\label{tab:contexts}
\end{table}

\textbf{Bối cảnh A (bệnh viện tuyến trung ương).} Ràng buộc không nằm ở hạ
tầng, mà ở việc tích hợp với một hệ thống PACS sẵn có và được sử dụng với cường
độ cao. Hệ thống không được làm xáo trộn quy trình làm việc đã định hình, phải
duy trì thông lượng lớn và phải tích hợp với phần mềm xem ảnh hiện hữu thay vì
thay thế nó. Lợi ích nằm ở việc bổ sung những năng lực mà hệ thống sẵn có còn
thiếu, chẳng hạn truy cập hình ảnh qua trình duyệt hay chia sẻ ca chụp giữa các
khoa, mà không làm thay đổi quy trình đang vận hành.

\textbf{Bối cảnh B (bệnh viện tuyến tỉnh).} Có hệ thống PACS nhưng có thể là
sản phẩm đời cũ với mức hỗ trợ DICOMweb hạn chế. Có một phần năng lực máy chủ;
việc trang bị GPU thì chưa chắc chắn. Bác sĩ chẩn đoán hình ảnh có mặt trong
giờ hành chính, nên khoảng thời gian ngoài giờ hành chính là nơi hệ thống đóng
góp nhiều nhất. Đây là trường hợp phổ biến nhất và là mục tiêu chính.

\textbf{Bối cảnh C (bệnh viện tuyến huyện).} Có máy CT nhưng không có PACS,
hình ảnh nằm trên máy trạm của thiết bị chụp. Không có phòng máy chủ, không có
nhân sự CNTT và không có bác sĩ chẩn đoán hình ảnh tại chỗ; kết nối tới cơ sở
tuyến trên có thể không liên tục. Ở đây hệ thống phải đảm nhiệm chính vai trò
lưu trữ, phải tự cài đặt và tự bảo trì với can thiệp tối thiểu, và phải hoạt
động được cả khi mất kết nối mạng. Vai trò chính của nó là bảo đảm hình ảnh
được lưu trữ an toàn ngay tại chỗ và chuyển được tới bác sĩ ở cơ sở tuyến trên
khi cần, khiến liên kết chẩn đoán hình ảnh từ xa trở thành thành phần thiết yếu
chứ không phải tuỳ chọn.

\subsection{Các yêu cầu phi chức năng xuyên suốt}

Bất kể bối cảnh nào, những điều sau đây đều đúng. \emph{Sự cố phải an toàn}:
nếu các thành phần xử lý phụ trợ không khả dụng, hình ảnh vẫn phải tiếp tục
được chuyển tới PACS và phần mềm xem ảnh mà không bị cản trở, do đó các thành
phần này không bao giờ được nằm trên đường tới hạn của việc truyền tải hình
ảnh. \emph{Dữ liệu phải truy được nguồn gốc}, kèm theo dấu thời gian và thông
tin về nơi phát sinh, để một ca chụp trong quá khứ luôn có thể được đối chiếu
lại. \emph{Dữ liệu đã tiếp nhận không được mất}: khi hệ thống đã xác nhận nhận ảnh từ
thiết bị chụp, bản duy nhất của ca chụp có thể chỉ còn nằm trong hệ thống, nên yêu
cầu độ bền áp dụng ở mọi bối cảnh, chỉ khác nhau ở cách hiện thực
(Mục~\ref{sec:durability}). \emph{Hệ thống phải cấu hình được khi triển khai}, từ quy tắc định tuyến
và thời hạn lưu trữ cho tới phân quyền truy cập, bởi mỗi bệnh viện có quy trình
làm việc riêng.
